\documentclass[aspectratio=169,professionalfonts]{beamer}

% Compile from this directory with:
% latexmk -pdf -interaction=nonstopmode -halt-on-error -outdir=build "presentation.tex"

\usepackage[T1]{fontenc}
\usepackage[utf8]{inputenc}
\usepackage{amsmath,amssymb,mathtools}
\usepackage{booktabs}
\usepackage{array}

\pdfstringdefDisableCommands{%
  \def\translate#1{#1}%
}

% ---------- Presentation metadata ----------
\title[Where Basic BRST Fails]{Where Basic BRST Becomes Insufficient}
\subtitle{Reducibility, open algebras, and the BV master equation}
\author[Group Presentation]{Group Presentation}
\institute{Gauge Field Theory by Ying Chen \\ Spring 2026}
\date{\today}

% ---------- Visual system ----------
\definecolor{CourseInk}{HTML}{1D2433}
\definecolor{CourseBlue}{HTML}{2454A6}
\definecolor{CourseTeal}{HTML}{16827A}
\definecolor{CourseGold}{HTML}{B78628}
\definecolor{CourseRed}{HTML}{B54750}
\definecolor{CourseMist}{HTML}{F3F5F7}
\definecolor{CourseLine}{HTML}{CED4DA}

\mode<presentation>{
  \useinnertheme{rectangles}
  \useoutertheme{default}
  \usefonttheme{professionalfonts}
}

\setbeamercolor{normal text}{fg=CourseInk,bg=white}
\setbeamercolor{structure}{fg=CourseBlue}
\setbeamercolor{frametitle}{fg=CourseInk,bg=white}
\setbeamercolor{framesubtitle}{fg=CourseTeal,bg=white}
\setbeamercolor{block title}{fg=white,bg=CourseBlue}
\setbeamercolor{block body}{fg=CourseInk,bg=CourseMist}
\setbeamercolor{block title alerted}{fg=white,bg=CourseRed}
\setbeamercolor{block body alerted}{fg=CourseInk,bg=CourseMist}
\setbeamercolor{block title example}{fg=white,bg=CourseTeal}
\setbeamercolor{block body example}{fg=CourseInk,bg=CourseMist}
\setbeamercolor{foot left}{fg=CourseInk,bg=CourseMist}
\setbeamercolor{foot right}{fg=white,bg=CourseBlue}
\setbeamercolor{course band}{fg=CourseInk,bg=CourseMist}

\setbeamerfont{title}{series=\bfseries,size=\LARGE}
\setbeamerfont{subtitle}{size=\large}
\setbeamerfont{frametitle}{series=\bfseries,size=\large}
\setbeamerfont{framesubtitle}{size=\small}
\setbeamerfont{block title}{series=\bfseries}
\setbeamerfont{footline}{size=\scriptsize}

\setbeamertemplate{navigation symbols}{}
\setbeamertemplate{itemize item}{\color{CourseTeal}\small$\blacktriangleright$}
\setbeamertemplate{itemize subitem}{\color{CourseGold}\scriptsize$\bullet$}
\setbeamertemplate{enumerate item}{\color{CourseBlue}\insertenumlabel.}
\setbeamertemplate{blocks}[rounded][shadow=false]

\setbeamertemplate{frametitle}{%
  \vspace{0.28cm}%
  \hspace*{0.45cm}\begin{minipage}{0.84\paperwidth}
    \usebeamerfont{frametitle}\insertframetitle\par
    \ifx\insertframesubtitle\empty
    \else
      \vspace{0.05cm}{\usebeamerfont{framesubtitle}\usebeamercolor[fg]{framesubtitle}\insertframesubtitle\par}
    \fi
    \vspace{0.10cm}{\color{CourseLine}\rule{\linewidth}{0.6pt}}
  \end{minipage}
  \vspace{-0.05cm}
}

\setbeamertemplate{footline}{%
  \leavevmode%
  \hbox{%
    \begin{beamercolorbox}[wd=.74\paperwidth,ht=2.6ex,dp=1.1ex,leftskip=0.45cm]{foot left}%
      \usebeamerfont{footline}\insertshorttitle
    \end{beamercolorbox}%
    \begin{beamercolorbox}[wd=.26\paperwidth,ht=2.6ex,dp=1.1ex,rightskip=0.45cm plus1fil]{foot right}%
      \usebeamerfont{footline}\hfill\insertframenumber/\inserttotalframenumber
    \end{beamercolorbox}%
  }%
}

\setbeamertemplate{title page}{%
  \vspace*{0.75cm}
  {\color{CourseBlue}\rule{0.18\paperwidth}{2.2pt}\par}
  \vspace{0.45cm}
  {\usebeamerfont{title}\usebeamercolor[fg]{normal text}\inserttitle\par}
  \vspace{0.18cm}
  {\usebeamerfont{subtitle}\color{CourseTeal}\insertsubtitle\par}
  \vspace{0.70cm}
  {\usebeamerfont{author}\insertauthor\par}
  \vspace{0.12cm}
  {\usebeamerfont{institute}\insertinstitute\par}
  \vfill
  \begin{beamercolorbox}[wd=\linewidth,sep=0.28cm,rounded=true]{course band}
    \small Gauge consistency as one equation on the extended field space.
    \hfill \insertdate
  \end{beamercolorbox}
}

\AtBeginSection[]{
  \begin{frame}[plain,noframenumbering]
    \vfill
    {\color{CourseBlue}\rule{0.16\paperwidth}{2pt}\par}
    \vspace{0.35cm}
    {\LARGE\bfseries\insertsection\par}
    \vspace{0.12cm}
    {\color{CourseTeal}Gauge field theory group discussion}
    \vfill
  \end{frame}
}

% ---------- Small helper commands ----------
\newcommand{\keyterm}[1]{\textcolor{CourseBlue}{\textbf{#1}}}
\newcommand{\sidehead}[1]{\textcolor{CourseTeal}{\textbf{#1}}}
\newcommand{\slidecite}[1]{{\scriptsize\color{CourseGold}#1}}
\newcommand{\gh}{\ensuremath{\operatorname{gh}}}
\newcommand{\EOM}{\ensuremath{\operatorname{EOM}}}
\newcommand{\BV}{\ensuremath{S_{\mathrm{BV}}}}

\begin{document}

\begin{frame}[plain]
  \titlepage
\end{frame}

\begin{frame}{Talk Map}
  \tableofcontents
\end{frame}

\section{Motivation}

\begin{frame}{Thesis}
  \begin{alertblock}{Main claim}
    Basic BRST is cleanest when the gauge algebra is closed and irreducible.
    BV is the systematic extension that also handles reducible symmetries and
    gauge algebras that close only modulo equations of motion.
  \end{alertblock}

  \vspace{0.25cm}
  \begin{columns}[T,onlytextwidth]
    \column{0.31\textwidth}
    \begin{block}{Reducibility}
      Gauge parameters can have their own gauge redundancy.
    \end{block}

    \column{0.31\textwidth}
    \begin{block}{Open algebra}
      Commutators can close only up to \EOM{} terms.
    \end{block}

    \column{0.31\textwidth}
    \begin{block}{BV answer}
      Add ghosts, antifields, and solve one master equation.
    \end{block}
  \end{columns}
\end{frame}

\begin{frame}{The Two Obstructions}
  \begin{columns}[T,onlytextwidth]
    \column{0.47\textwidth}
    \begin{exampleblock}{2-form gauge theory}
      \[
        \Lambda_\mu \sim \Lambda_\mu+\partial_\mu\alpha
      \]
      The gauge parameter is itself redundant, so a single vector ghost leaves
      a zero mode in the Faddeev-Popov operator.
    \end{exampleblock}

    \column{0.47\textwidth}
    \begin{exampleblock}{Poisson sigma model}
      \[
        [\delta_{\epsilon_1},\delta_{\epsilon_2}]
        =\delta_{\epsilon_3}+\EOM
      \]
      The algebra knows about equations of motion, so off-shell nilpotency
      cannot be imposed on the original fields alone.
    \end{exampleblock}
  \end{columns}
\end{frame}

\begin{frame}{2-Form: Reducible Gauge Parameter}
  \sidehead{Free Abelian 2-form}
  \[
    S_0[B]=-\frac{1}{12}\int d^Dx\,H_{\mu\nu\rho}H^{\mu\nu\rho},
    \qquad
    H_{\mu\nu\rho}=3\partial_{[\mu}B_{\nu\rho]} .
  \]

  \sidehead{Gauge symmetry}
  \[
    \delta_\Lambda B_{\mu\nu}
    =\partial_\mu\Lambda_\nu-\partial_\nu\Lambda_\mu .
  \]

  \begin{block}{Parameter redundancy}
    \[
      \Lambda_\mu\sim\Lambda_\mu+\partial_\mu\alpha,
      \qquad
      \delta_{\Lambda=d\alpha}B_{\mu\nu}=0 .
    \]
  \end{block}
\end{frame}

\begin{frame}{2-Form: Naive Faddeev-Popov Operator}
  Choose the Lorenz-type gauge condition
  \[
    G_\nu(B)=\partial^\mu B_{\mu\nu}=0 .
  \]

  Its variation is
  \[
    \delta_\Lambda G_\nu=M_\nu{}^\rho\Lambda_\rho,
    \qquad
    M_\nu{}^\rho=\Box\delta_\nu{}^\rho-\partial_\nu\partial^\rho .
  \]

  The naive gauge-fixed expression would contain
  \[
    Z_{\mathrm{naive}}
    =\int DB\,\delta[G(B)]\det M\,e^{iS_0[B]} .
  \]
\end{frame}

\begin{frame}{2-Form: One Vector Ghost Is Not Enough}
  The vector ghost action would be
  \[
    S_{\mathrm{gh}}
    =\int d^Dx\,\bar c^\nu M_\nu{}^\rho c_\rho .
  \]

  But exact gauge parameters lie in the kernel:
  \[
    M_\nu{}^\rho\partial_\rho\alpha
    =\Box\partial_\nu\alpha-\partial_\nu\Box\alpha=0 .
  \]

  \begin{alertblock}{Consequence}
    The determinant is singular because the gauge fixing did not divide by
    the redundancy of the gauge parameter.
  \end{alertblock}
\end{frame}

\begin{frame}{Poisson Sigma Model: Data}
  \small
  \begin{columns}[T,onlytextwidth]
    \column{0.45\textwidth}
    \sidehead{Fields}
    \[
      \begin{aligned}
        X^i&:\Sigma\to M,\\
        A_i&\in\Omega^1(\Sigma,X^*T^*M).
      \end{aligned}
    \]

    \sidehead{Action}
    \[
      \begin{aligned}
        S_0[X,A]
        &=\int_\Sigma A_i\wedge dX^i\\
        &\quad+\frac12\pi^{ij}(X)A_i\wedge A_j .
      \end{aligned}
    \]

    \column{0.49\textwidth}
    \sidehead{Gauge transformations}
    \[
      \delta_\epsilon X^i=\pi^{ij}(X)\epsilon_j ,
    \]
    \[
      \begin{aligned}
        \delta_\epsilon A_i
        &=-d\epsilon_i\\
        &\quad-\partial_i\pi^{jk}(X)A_j\epsilon_k .
      \end{aligned}
    \]
  \end{columns}
\end{frame}

\begin{frame}{Poisson Sigma Model: On-Shell Closure}
  For a general Poisson tensor,
  \[
    [\delta_{\epsilon_1},\delta_{\epsilon_2}]
    =\delta_{\epsilon_3}+\EOM,
    \qquad
    \epsilon_{3i}
    =\partial_i\pi^{jk}\epsilon_{1j}\epsilon_{2k}.
  \]

  The equations of motion are
  \[
    dX^i+\pi^{ij}A_j=0,
    \qquad
    dA_i+\frac12\partial_i\pi^{jk}A_j\wedge A_k=0 .
  \]

  \begin{block}{Message}
    Nilpotency cannot be made off-shell without recording the \EOM{} part of
    the algebra.
  \end{block}
\end{frame}

\begin{frame}{Path Integral Obstruction}
  Gauge fixing by a BRST-exact deformation would suggest
  \[
    Z_\Psi=\int D\Phi\,
    \exp\left[\frac{i}{\hbar}(S_0+s\Psi)\right].
  \]

  For an open algebra,
  \[
    s(S_0+s\Psi)=s^2\Psi
    \sim
    \frac{\delta\Psi}{\delta\Phi^I}
    N^{IJ}\frac{\delta S_0}{\delta\Phi^J}.
  \]

  \[
    \left\langle F\frac{\delta S_0}{\delta\Phi^I}\right\rangle
    =i\hbar\left\langle\frac{\delta F}{\delta\Phi^I}\right\rangle .
  \]

  \begin{alertblock}{Interpretation}
    \EOM{} insertions are not zero as operator insertions in the path integral.
  \end{alertblock}
\end{frame}

\begin{frame}{Moral}
  \begin{itemize}
    \item Reducible theories need ghosts for the redundancies of gauge
      parameters.
    \item Open algebras need a formalism that remembers which terms are
      proportional to equations of motion.
    \item BV does both by enlarging the field space and solving one equation.
  \end{itemize}

  \vspace{0.25cm}
  \begin{block}{BV slogan}
    Gauge transformations, their algebra, reducibility, Jacobi identities, and
    \EOM{} terms are all encoded in the master action.
  \end{block}
\end{frame}

\section{BV Construction}

\begin{frame}{General Input}
  Start from a gauge theory
  \[
    S_0[\phi],
    \qquad
    \delta_\epsilon\phi^i=R^i_\alpha(\phi)\epsilon^\alpha .
  \]

  Gauge invariance means
  \[
    \frac{\delta S_0}{\delta\phi^i}R^i_\alpha=0 .
  \]

  Allow both complications:
  \[
    R^i_\alpha Z^\alpha_a=0,
    \qquad
    [R_\alpha,R_\beta]^i
    =R^i_\gamma f^\gamma_{\alpha\beta}
    +\frac{\delta S_0}{\delta\phi^j}M^{ji}_{\alpha\beta}.
  \]
\end{frame}

\begin{frame}{Fields, Ghosts, Antifields}
  Add ghosts for every stage of the gauge complex:
  \[
    \Phi^A=(\phi^i,c^\alpha,c^a_{(1)},\ldots).
  \]

  Add an antifield for every field:
  \[
    \Phi_A^*=(\phi_i^*,c_\alpha^*,c^*_{(1)a},\ldots),
    \qquad
    \gh(\Phi_A^*)=-1-\gh(\Phi^A).
  \]

  \begin{exampleblock}{Role of antifields}
    Antifields are not new propagating particles. They are sources that record
    transformations, algebra, reducibility, and \EOM{} terms.
  \end{exampleblock}
\end{frame}

\begin{frame}{BV Antibracket}
  The BV antibracket pairs each field with its antifield:
  \[
    (F,G)
    =\int
    \left(
    \frac{\delta_R F}{\delta\Phi^A}
    \frac{\delta_L G}{\delta\Phi_A^*}
    -
    \frac{\delta_R F}{\delta\Phi_A^*}
    \frac{\delta_L G}{\delta\Phi^A}
    \right).
  \]

  The BV differential is Hamiltonian:
  \[
    s_{\mathrm{BV}}F=(F,\BV).
  \]

  \begin{block}{Useful picture}
    The antibracket turns the master action into the generator of all BRST/BV
    transformations on the extended space.
  \end{block}
\end{frame}

\begin{frame}{Minimal BV Action: First Layers}
  The minimal action starts with the classical action and adds the gauge data:
  \[
    S_{\min}
    =S_0+\phi_i^*R^i_\alpha c^\alpha+\cdots .
  \]

  The next terms record the possible complications:
  \[
    S_{\min}\supset
    c_\alpha^*Z^\alpha_a c^a_{(1)}
    +c_\gamma^*f^\gamma_{\alpha\beta}c^\alpha c^\beta
    +\phi_i^*\phi_j^*M^{ji}_{\alpha\beta}c^\alpha c^\beta .
  \]

  \slidecite{Precise coefficients and signs are fixed by the master equation and by parity conventions.}
\end{frame}

\begin{frame}{What the Antifield Terms Mean}
  \small
  \begin{tabular}{>{\raggedright\arraybackslash}p{0.30\textwidth}
                  >{\raggedright\arraybackslash}p{0.35\textwidth}
                  >{\raggedright\arraybackslash}p{0.24\textwidth}}
    \toprule
    Term & Encoded datum & Why it matters \\
    \midrule
    $\phi_i^*R^i_\alpha c^\alpha$ &
    Gauge transformations &
    Generates $s\phi^i$. \\
    $c_\alpha^*Z^\alpha_a c^a_{(1)}$ &
    Reducibility &
    Adds ghost-for-ghosts. \\
    $c_\gamma^*f^\gamma_{\alpha\beta}c^\alpha c^\beta$ &
    Gauge algebra &
    Generates the ghost bracket. \\
    $\phi_i^*\phi_j^*M^{ji}_{\alpha\beta}c^\alpha c^\beta$ &
    Open-algebra \EOM{} term &
    Restores off-shell nilpotency. \\
    \bottomrule
  \end{tabular}
\end{frame}

\begin{frame}{The Master Equation}
  The central condition is
  \[
    (\BV,\BV)=0 .
  \]

  It simultaneously imposes:
  \begin{itemize}
    \item invariance of $S_0$;
    \item closure of gauge transformations;
    \item reducibility relations;
    \item Jacobi and higher identities;
    \item cancellation of open-algebra \EOM{} terms.
  \end{itemize}

  \begin{alertblock}{One equation}
    The master equation is the compact replacement for many separate
    consistency checks.
  \end{alertblock}
\end{frame}

\begin{frame}{Off-Shell Nilpotency}
  If the master equation holds, the BV differential squares to zero:
  \[
    s_{\mathrm{BV}}^2F
    =\frac12(F,(\BV,\BV))=0 .
  \]

  \begin{block}{Point of the extension}
    The nilpotency is off shell on the extended field-antifield space, even
    when the original gauge algebra closed only modulo \EOM{} terms.
  \end{block}

  \vspace{0.2cm}
  Closed irreducible BRST is the special case where there are no
  ghost-for-ghosts and no antifield-quadratic open-algebra corrections.
\end{frame}

\begin{frame}{Closed BRST as a Truncation}
  If the theory is closed and irreducible,
  \[
    Z=0,
    \qquad
    M=0 .
  \]

  Then
  \[
    S_{\min}
    =S_0
    +\phi_i^*R^i_\alpha c^\alpha
    +\frac12 c_\gamma^*f^\gamma_{\alpha\beta}c^\alpha c^\beta .
  \]

  \begin{exampleblock}{Interpretation}
    Familiar BRST is the easiest truncation of BV, not the most general
    starting point.
  \end{exampleblock}
\end{frame}

\section{Gauge Fixing}

\begin{frame}{Why Antifields Are Useful}
  Before gauge fixing, antifields expose the equations of motion:
  \[
    s_{\mathrm{BV}}\phi_i^*
    =\frac{\delta S_0}{\delta\phi^i}+\cdots .
  \]

  After gauge fixing, they are eliminated by a choice of Lagrangian submanifold:
  \[
    \Phi_A^*=\frac{\delta\Psi}{\delta\Phi^A}.
  \]

  \begin{block}{Practical meaning}
    Antifields are sources for consistency before gauge fixing, then become
    derivatives of the gauge-fixing fermion.
  \end{block}
\end{frame}

\begin{frame}{Non-Minimal Pair}
  Minimal BV knows the gauge symmetry, but it has not chosen a gauge.

  To impose
  \[
    G^\alpha(\phi)=0,
  \]
  add a contractible pair
  \[
    s\bar c_\alpha=b_\alpha,
    \qquad
    sb_\alpha=0 .
  \]

  In the BV action:
  \[
    S_{\mathrm{nm}}=S_{\min}+\int \bar c^{*\alpha}b_\alpha .
  \]
\end{frame}

\begin{frame}{Gauge-Fixing Fermion}
  The gauge-fixing fermion has ghost number $-1$.

  For a gauge condition $G^\alpha(\phi)=0$, choose
  \[
    \Psi
    =\int \bar c_\alpha
    \left(G^\alpha(\phi)+\frac{\xi}{2}b^\alpha\right).
  \]

  Then set
  \[
    \bar c^{*\alpha}=G^\alpha+\frac{\xi}{2}b^\alpha,
    \qquad
    \phi_i^*=\bar c_\alpha\frac{\delta G^\alpha}{\delta\phi^i}.
  \]
\end{frame}

\begin{frame}{What Gauge Fixing Produces}
  After substituting $\Phi^*=\delta\Psi/\delta\Phi$,
  \[
    S_\Psi
    =S_0
    +b_\alpha G^\alpha
    +\frac{\xi}{2}b_\alpha b^\alpha
    -\bar c_\alpha M^\alpha{}_\beta c^\beta .
  \]

  Here
  \[
    M^\alpha{}_\beta
    =\frac{\delta G^\alpha}{\delta\phi^i}R^i_\beta
  \]
  is the Faddeev-Popov operator.

  \begin{block}{Roles}
    $b_\alpha$ imposes the gauge condition, while
    $\bar c_\alpha,c^\alpha$ produce the FP determinant.
  \end{block}
\end{frame}

\begin{frame}{BV Path Integral}
  The gauge-fixed action is
  \[
    S_\Psi(\Phi)
    =
    \BV\left(\Phi,\Phi^*=\frac{\delta\Psi}{\delta\Phi}\right).
  \]

  Then
  \[
    Z_\Psi
    =\int D\Phi\,
    \exp\left[\frac{i}{\hbar}S_\Psi(\Phi)\right].
  \]

  At quantum level the classical master equation becomes
  \[
    \frac12(\BV,\BV)-i\hbar\Delta\BV=0 .
  \]
\end{frame}

\section{Worked Examples}

\begin{frame}{2-Form BV Field Content}
  Minimal fields:
  \[
    B_{\mu\nu},
    \qquad
    c_\mu,
    \qquad
    \rho .
  \]

  Meanings:
  \begin{itemize}
    \item $c_\mu$ is the ghost for $\Lambda_\mu$;
    \item $\rho$ is the ghost for $\Lambda_\mu\sim\Lambda_\mu+\partial_\mu\alpha$.
  \end{itemize}

  Antifields:
  \[
    B^{*\mu\nu},
    \qquad
    c^{*\mu},
    \qquad
    \rho^* .
  \]
\end{frame}

\begin{frame}{2-Form Minimal BV Action}
  \[
    S_{\min}
    =S_0[B]
    +\int d^Dx\,B^{*\mu\nu}
    (\partial_\mu c_\nu-\partial_\nu c_\mu)
    +\int d^Dx\,c^{*\mu}\partial_\mu\rho .
  \]

  This gives
  \[
    sB_{\mu\nu}=\partial_\mu c_\nu-\partial_\nu c_\mu,
    \qquad
    sc_\mu=\partial_\mu\rho,
    \qquad
    s\rho=0 .
  \]
\end{frame}

\begin{frame}{2-Form Nilpotency Check}
  Now the missing ghost direction is included:
  \[
    s^2B_{\mu\nu}
    =\partial_\mu\partial_\nu\rho
    -\partial_\nu\partial_\mu\rho=0 .
  \]

  \begin{exampleblock}{Complex}
    \[
      \rho
      \xrightarrow{\ d\ }
      c_\mu
      \xrightarrow{\ d\ }
      B_{\mu\nu}.
    \]
  \end{exampleblock}

  The naive zero mode becomes a genuine part of the BV complex.
\end{frame}

\begin{frame}{2-Form Gauge Fixing}
  A schematic gauge-fixing fermion is
  \[
    \Psi
    =\int d^Dx\,
    \bar c^\nu\partial^\mu B_{\mu\nu}
    +\int d^Dx\,\bar\rho\,\partial^\mu c_\mu+\cdots .
  \]

  \begin{itemize}
    \item $\bar c^\nu\partial^\mu B_{\mu\nu}$ fixes the 2-form gauge symmetry.
    \item $\bar\rho\,\partial^\mu c_\mu$ fixes the exact direction of the vector ghost.
  \end{itemize}

  \begin{block}{Message}
    BV fixes the whole complex, not just the original gauge field.
  \end{block}
\end{frame}

\begin{frame}{2-Form Determinant Cure}
  The naive operator has zero modes:
  \[
    M_1=d^\dagger d,
    \qquad
    \ker M_1\supset \operatorname{im}d .
  \]

  BV gauge fixing also imposes
  \[
    d^\dagger B=0,
    \qquad
    d^\dagger c=0 .
  \]

  \[
    \det M_1
    \quad\leadsto\quad
    \det{}'\!\left(d^\dagger d\big|_{\Omega^1/\operatorname{im}d}\right)
    \times
    \text{scalar ghost factor}.
  \]
\end{frame}

\begin{frame}{Poisson Sigma Model BV Fields}
  Minimal fields:
  \[
    X^i,
    \qquad
    A_i,
    \qquad
    c_i .
  \]

  The ghost $c_i$ replaces the gauge parameter $\epsilon_i$.

  Antifields:
  \[
    X_i^*,
    \qquad
    A^{*i},
    \qquad
    c^{*i}.
  \]

  \begin{block}{Open-algebra expectation}
    Because the gauge algebra closes only on shell, the BV action should
    contain an antifield-quadratic correction.
  \end{block}
\end{frame}

\begin{frame}{Poisson Sigma Model BV Action}
  \small
  One common component form, up to sign conventions, is
  \[
  \begin{aligned}
    S_{\min}=S_0
    &+\int_\Sigma X_i^*\pi^{ij}c_j
    -\int_\Sigma A^{*i}\wedge
    \left(dc_i+\partial_i\pi^{jk}A_jc_k\right)\\
    &+\frac12\int_\Sigma c^{*i}\partial_i\pi^{jk}c_jc_k
    +\frac14\int_\Sigma A^{*i}\wedge A^{*j}
    \partial_i\partial_j\pi^{kl}c_kc_l .
  \end{aligned}
  \]

  \begin{alertblock}{Key term}
    The $A^*A^*$ term is the BV correction that records the open-algebra
    \EOM{} contribution.
  \end{alertblock}
\end{frame}

\begin{frame}{Reading the Poisson Sigma Terms}
  At antifields set to zero:
  \[
    sX^i=\pi^{ij}c_j .
  \]

  The one-form transforms as
  \[
    sA_i=-dc_i-\partial_i\pi^{jk}A_jc_k+\cdots .
  \]

  The ghost algebra is encoded by
  \[
    sc_i\sim \partial_i\pi^{jk}c_jc_k .
  \]

  The antifield-quadratic term records the \EOM{} part of the closure.
\end{frame}

\begin{frame}{Why the Correction Works}
  The master equation uses the Poisson identity
  \[
    \pi^{\ell[i}\partial_\ell\pi^{jk]}=0 .
  \]

  Without the $A^*A^*$ term, the master equation would leave uncanceled
  \EOM{} terms. With it,
  \[
    (S_{\min},S_{\min})=0 .
  \]

  \begin{exampleblock}{Conclusion}
    The BRST differential becomes off-shell nilpotent only on the full BV
    field-antifield space.
  \end{exampleblock}
\end{frame}

\section{String Field Theory}

\begin{frame}{Application: Cubic Open String Field Theory}
  Witten's cubic open string field theory has schematic action
  \[
    S_{\mathrm{OSFT}}
    =\frac12\langle \Psi,Q_B\Psi\rangle
    +\frac{g_o}{3}\langle \Psi,\Psi*\Psi\rangle .
  \]

  \begin{itemize}
    \item $\Psi$ is the open string field.
    \item $Q_B$ is the worldsheet BRST operator.
    \item $*$ is the open-string star product.
    \item $\langle-,-\rangle$ is the BPZ pairing.
  \end{itemize}

  \[
    \delta\Psi
    =Q_B\Lambda+g_o(\Psi*\Lambda-\Lambda*\Psi).
  \]
\end{frame}

\begin{frame}{Open String Field Theory as BV}
  In the usual classical action, $\Psi$ has ghost number $1$.

  In the BV formulation, $\Psi$ contains components of all ghost numbers.
  Since the BPZ pairing is nonzero when the total ghost number is $3$,
  components of ghost number $g$ pair with components of ghost number $3-g$.

  \begin{block}{BV interpretation}
    The same cubic functional becomes a BV master action after the string
    field is enlarged to include fields and antifields simultaneously.
  \end{block}
\end{frame}

\begin{frame}{Algebra Behind the OSFT Master Equation}
  The master equation follows from cyclicity of the BPZ pairing and
  \[
    Q_B^2=0,
  \]
  \[
    Q_B(A*B)=(Q_BA)*B+(-1)^{|A|}A*(Q_BB),
  \]
  \[
    (A*B)*C=A*(B*C).
  \]

  Thus consistency of the cubic vertex is encoded as
  \[
    (S_{\mathrm{OSFT}},S_{\mathrm{OSFT}})=0 .
  \]
\end{frame}

\section{Close}

\begin{frame}{Practical BV Recipe}
  BV is a construction:
  \begin{enumerate}
    \item Write the gauge transformations.
    \item Add ghosts and antifields.
    \item Add antifield-linear terms for transformations.
    \item Add higher ghost terms for reducibility.
    \item Add antifield-quadratic terms for open algebras.
    \item Solve $(S,S)=0$.
    \item Choose $\Psi$ and integrate over gauge-fixed fields.
  \end{enumerate}
\end{frame}

\begin{frame}{Closing Point}
  \begin{alertblock}{One controlling equation}
    BV turns gauge consistency into the master equation
    \[
      (S,S)=0 .
    \]
  \end{alertblock}

  \begin{itemize}
    \item Reducible theories require the ghost-for-ghost complex.
    \item Open algebras require antifield-dependent correction terms.
    \item Open string field theory packages the cubic vertex consistency in
      the same BV language.
  \end{itemize}
\end{frame}

\appendix

\section{Backup}

\begin{frame}{Backup: General Open Algebra}
  \small
  The general closure relation has the form
  \[
    [R_\alpha,R_\beta]^i
    =
    R^i_\gamma f^\gamma_{\alpha\beta}
    +
    \frac{\delta S_0}{\delta\phi^j}M^{ji}_{\alpha\beta}.
  \]

  The antifield-quadratic term in $S_{\min}$ is the local bookkeeping device
  for the second term:
  \[
    S_{\min}\supset
    \phi_i^*\phi_j^*M^{ji}_{\alpha\beta}c^\alpha c^\beta .
  \]

  This is why open algebras naturally push us beyond the field-only BRST
  complex.
\end{frame}

\begin{frame}{Backup: Gauge-Fixing Substitution}
  \small
  For
  \[
    \Psi
    =
    \int \bar c_\alpha
    \left(G^\alpha(\phi)+\frac{\xi}{2}b^\alpha\right),
  \]
  the relevant antifields are
  \[
    \bar c^{*\alpha}
    =G^\alpha+\frac{\xi}{2}b^\alpha,
    \qquad
    \phi_i^*
    =\bar c_\alpha\frac{\delta G^\alpha}{\delta\phi^i}.
  \]

  Substitution into the non-minimal action gives
  \[
    b_\alpha G^\alpha
    +\frac{\xi}{2}b_\alpha b^\alpha
    -\bar c_\alpha
    \frac{\delta G^\alpha}{\delta\phi^i}R^i_\beta c^\beta .
  \]
\end{frame}

\begin{frame}{Backup: 2-Form Complex}
  The reducible gauge structure is the differential complex
  \[
    \Omega^0
    \xrightarrow{\ d\ }
    \Omega^1
    \xrightarrow{\ d\ }
    \Omega^2 .
  \]

  In components:
  \[
    s\rho=0,
    \qquad
    sc_\mu=\partial_\mu\rho,
    \qquad
    sB_{\mu\nu}=\partial_\mu c_\nu-\partial_\nu c_\mu .
  \]

  The identity $d^2=0$ gives $s^2=0$ on the minimal complex.
\end{frame}

\begin{frame}{Backup: Poisson Sigma Conventions}
  \small
  The action and gauge transformations used here are
  \[
    S_0[X,A]
    =\int_\Sigma A_i\wedge dX^i
    +\frac12\pi^{ij}(X)A_i\wedge A_j ,
  \]
  \[
    \delta_\epsilon X^i=\pi^{ij}(X)\epsilon_j,
    \qquad
    \delta_\epsilon A_i
    =-d\epsilon_i-\partial_i\pi^{jk}(X)A_j\epsilon_k .
  \]

  Different sign conventions move signs between the $A^*$ and $c^*$ terms,
  but the structural role of the antifield-quadratic term is unchanged.
\end{frame}

\begin{frame}{Backup: OSFT Ghost Number Pairing}
  \small
  The BPZ pairing pairs components whose ghost numbers add to $3$:
  \[
    \langle \Psi_g,\Psi_{3-g}\rangle\neq 0 .
  \]

  Allowing the string field to contain all ghost numbers packages fields and
  antifields into one object:
  \[
    \Psi=\sum_g \Psi_g .
  \]

  The cubic OSFT functional is therefore naturally suited to the BV language.
\end{frame}

\end{document}
